% ===============================================================
\chapter*{Citations}
\addcontentsline{toc}{chapter}{Citations}
\markboth{CITATIONS}{CITATIONS}
% ===============================================================

\section*{Main source}

The slides, and most of the content of these notes, are taken from the course \emph{Introduction to Philosophy of Language} (PHI-2300/ENG-2910), taught by Professor Eric Snyder at Ashoka University in the Monsoon 2025 semester. This course is the main citation for these notes. 

\section*{Course textbooks}

\begin{itemize}
  \item Martinich, \emph{The Philosophy of Language}, 3rd edition, 1996.
  \item Miller, \emph{Philosophy of Language}, 2007.
  \item McGinn, \emph{Philosophy of Language}, 2016.
  \item Lycan, \emph{Philosophy of Language}, 2008.
  \item Soames, \emph{Philosophy of Language}, 2010.
\end{itemize}

Another source of material is the MIT course on the Philosophy of Language:\\
\url{https://ocw.mit.edu/courses/24-251-introduction-to-philosophy-of-language-fall-2011/}

\section*{Readings by week}

Supplemental, non-required readings are marked with an asterisk ($\ast$).

\newcommand{\req}{\item[--]}
\newcommand{\supp}{\item[$\ast$]}
\newenvironment{readings}{\begin{itemize}[leftmargin=2em,itemsep=0pt]}{\end{itemize}}

\paragraph{Week 1: Preliminaries}
\begin{readings}
  \req Syllabus
  \req `Introduction', pp.\ 3--19
  \supp pp.\ xvii--xviii
  \supp pp.\ 1--6
  \supp \url{https://iep.utm.edu/lang-phi/#H5}
\end{readings}

\paragraph{Week 2: Frege: Reference and Compositionality}
\begin{readings}
  \req Gottlob Frege, `On Sense and Nominatum', pp.\ 186--189
  \req `Frege: Semantic Value and Reference', pp.\ 1--21
  \supp pp.\ 1--6
  \supp pp.\ 1--17
  \supp pp.\ 7--13
  \supp \url{https://plato.stanford.edu/entries/reference/}
  \supp \url{https://plato.stanford.edu/entries/compositionality/}
\end{readings}

\paragraph{Week 3: Frege: Sense and Substitutivity}
\begin{readings}
  \req `Frege and Russell: Sense and Definite Descriptions', pp.\ 23--36, 43--52
  \req `Frege on Sense and Reference', pp.\ 1--24
  \supp pp.\ 31--37
  \supp pp.\ 13--20
  \supp \url{https://plato.stanford.edu/entries/frege/}
\end{readings}

\paragraph{Week 4: Russell's Theory of Descriptions}
\begin{readings}
  \req Bertrand Russell, `Descriptions', pp.\ 208--213
  \req `Russell on Definite Descriptions', pp.\ 55--72
  \supp pp.\ 64--69
  \supp pp.\ 9--18
  \supp pp.\ 20--26
  \supp \url{https://plato.stanford.edu/entries/descriptions/}
\end{readings}

\paragraph{Week 5: The Description Theory of Names}
\begin{readings}
  \req `Proper Names: the Description Theory', pp.\ 31--43
  \req `Frege and Russell: Sense and Definite Descriptions', pp.\ 64--77
  \supp pp.\ 26--32
  \supp Stephen Neale (1993), \emph{Descriptions}
  \supp \url{https://plato.stanford.edu/entries/descriptions/}
  \supp \url{https://plato.stanford.edu/entries/reference/}
\end{readings}

\paragraph{Week 6: Strawson on Referring}
\begin{readings}
  \req P.F. Strawson, `On Referring', pp.\ 215--224
  \req `Definite Descriptions', pp.\ 19--23
  \supp pp.\ 77--78
  \supp pp.\ 80--83
  \supp Bertrand Russell, `Mr. Strawson on Referring'
  \supp \url{https://plato.stanford.edu/entries/descriptions/}
\end{readings}

\paragraph{Week 7: Grice's Theory of Implicature}
\begin{readings}
  \req Paul Grice, `Logic and Conversation', pp.\ 156--163
  \req `Implicative Relations', pp.\ 86--96
  \supp pp.\ 246--253, 259--264
  \supp pp.\ 191--202
  \supp \url{https://plato.stanford.edu/entries/grice/}
  \supp \url{https://plato.stanford.edu/entries/implicature/}
\end{readings}

\paragraph{Week 8: Verificationism and Quine's Skepticism}
\begin{readings}
  \req W.V.O. Quine, `Two Dogmas of Empiricism', Intro and Sections 1, 5 and 6
  \req `Skepticism about Sense (I)', pp.\ 126--140
  \supp pp.\ 98--107
  \supp \url{https://plato.stanford.edu/entries/quine/}, \S3
  \supp \url{https://plato.stanford.edu/entries/logical-empiricism/#EmpVerAntMet}, \S4.1--2
\end{readings}

\paragraph{Week 9: Tarski and Davidson}
\begin{readings}
  \req Alfred Tarski, `The Semantic Conception of Truth', \S\S1--12 (pp.\ 61--70)
  \req `Sense and Truth: Tarski and Davidson', pp.\ 271--305 (section 8.5 is optional)
  \supp pp.\ 165--190
  \supp pp.\ 109--124
  \supp pp.\ 33--48
  \supp \url{https://plato.stanford.edu/entries/tarski-truth/}
  \supp \url{https://plato.stanford.edu/entries/davidson/}, \S3
\end{readings}

\paragraph{Week 10: Possible Worlds and Direct Reference}
\begin{readings}
  \req Saul Kripke, `Naming and Necessity', pp.\ 259--266
  \req `Kripke on Names', pp.\ 35--49
  \supp pp.\ 45--58
  \supp pp.\ 77--92
  \supp Saul Kripke (1980), \emph{Naming and Necessity}
  \supp \url{https://plato.stanford.edu/entries/rigid-designators/}
  \supp \url{https://plato.stanford.edu/entries/logic-modal-origins/}, \S3.2
\end{readings}

\paragraph{Week 11: Twin Earth and Semantic Externalism}
\begin{readings}
  \req Hilary Putnam, `Meaning and Reference', pp.\ 284--288
  \req `Putnam on Semantic Externalism', pp.\ 133--146
  \supp pp.\ 58--60
  \supp pp.\ 323--329
  \supp pp.\ 88--91
  \supp \url{https://plato.stanford.edu/entries/natural-kinds/}, \S4
  \supp \url{https://plato.stanford.edu/entries/skepticism-content-externalism/}, \S3
\end{readings}

\paragraph{Week 12: Kripke's Puzzle}
\begin{readings}
  \req Saul Kripke, `A Puzzle About Belief', pp.\ 392--395, 398--401
  \req `Kripke's Puzzle About Belief', pp.\ 203--209
  \supp \url{https://plato.stanford.edu/entries/prop-attitude-reports/}
\end{readings}
